\newpage
\part{General Information}

\textbf{Model Context Protocol (MCP)} is a standardized communication \textbf{protocol}
that securely connects AI large language \textbf{models} (LLMs) to external
\textbf{contexts} like codebases, databases, and files, frequently using either
standardized input and output (\textbf{STDIO}) for local integration or
Server-Sent Events (\textbf{HTTP/SSE}) for network-based remote communication.

Regardless of the transport, all MCP clients and servers exchange messages using
\textbf{JSON-RPC 2.0}: structured request, response, and notification payloads
sent over the chosen channel.

\section{MCP Hub}

People use \textbf{MCP Hub} in two different senses. The distinction matters when you read docs or community posts.

\subsection{Runtime Hub}

On your machine, the host or \textbf{runtime hub} inlcudes the following:
\begin{enumerate}
    \item Claude Desktop (CLI)
    \item Cursor (IDE)
    \item Windsurf (IDE)
    \item any another CLI or IDE AI agents
\end{enumerate} 
The \textbf{host} reads one configuration file, \uline{starts each MCP server as a background process}, 
and keeps every client session alive at the same time.
When the \textbf{model} calls a \textbf{tool}, the host routes the \textbf{JSON-RPC} request to the right server and passes the result back into the conversation.
In short, the host is the local \textbf{orchestrator} that wires many servers into a single agent session.

\subsection{Discovery Hub}

The same label is also used for a \textbf{public registry} or marketplace, like npm or Docker Hub, where you browse and install pre-built servers.
Sites such as \url{https://smithery.ai} and \url{https://mcp-get.com} catalog community integrations (Slack, GitHub, memory stores, and more), often with a one-command install.
Discovery hub \uline{helps you find servers, not run them}. Execution still happens in your host.
